\SourceLeaf{392}%
als Differentialgleichung sowohl für $p$ als auch für $q$ und für $r$, das ist für
$\surd{\frac{a}{\mathrm M(a,b)}}$, $\surd{\frac{b}{\mathrm M(a,b)}}$ und
$\surd{\frac{c}{\mathrm M(a,b)}}$ und ebenso auch, wenn man
$-\pi\frac{\mathrm M(a,c)}{\mathrm M(a,b)}=\log z$ statt $\log y$ darin gesetzt denkt,
als Differentialgleichung für $\surd{\frac{a}{\mathrm M(a,c)}}$,
$\surd{\frac{c}{\mathrm M(a,c)}}$ und $\surd{\frac{b}{\mathrm M(a,c)}}$.]

\AGMArticle{16}
[Die Darstellung der Quotienten der Grössen $a,b,c,\mathrm M(a,b),\mathrm M(a,c)$ durch
Reihen, die nach Potenzen von $y^{\frac14}$ oder $z^{\frac14}$ fortschreiten, lässt sich mit Hülfe
der Fundamentalsätze des hier zu untersuchenden Algorithmus z.\,B. in folgender
Weise ausführen.

Nach der Definition von $y$ und dem in Art.\,12 für den rückwärts verlängerten Algorithmus aufgestellten Satze wird:
\[
-\pi\frac{\mathrm M(a,b)}{\mathrm M(a,c)}
=-\pi\frac{\mathrm M(a^{\mathrm n},b^{\mathrm n})}{2^n\mathrm M(a^{\mathrm n},c^{\mathrm n})}=\log y
\]
also nach den Gleichungen in Art.\,12 der Grenzwerth von
$\frac12y^{-2^{n-2}}\cdot\surd{\frac{c^{\mathrm n}}{\mathrm M(a^{\mathrm n},b^{\mathrm n})}}$ für ein
immer wachsendes $n$, oder was dasselbe ist, der Grenzwerth von $\frac12y^{-\frac14}r(y)$,
wo $r(y)$ statt $\surd{\frac{c}{\mathrm M(a,b)}}$ gesetzt ist, für bis zur Null abnehmendes positives $y^{\frac14}$
gleich der Einheit.

Bezeichnen wir noch $\surd{\frac{a}{\mathrm M(a,b)}}$ und $\surd{\frac{b}{\mathrm M(a,b)}}$ durch $p(y)$ und $q(y)$, so folgt
aus $aa=bb+cc$ und den beiden Gleichungen für $\surd{\mathrm M(a,b)}$ in Art.\,12
\[
\begin{aligned}p(y)&=1+r(y^4)+r(y^{16})+r(y^{64})\\&\quad+.\;.+r(y^{2^{2n}})+.\;.\end{aligned}
\]
\[
\begin{aligned}q(y)&=1-r(y^4)+r(y^{16})+r(y^{64})\\&\quad+.\;.+r(y^{2^{2n}})+.\;.\end{aligned}
\]
\[
\begin{aligned}r(y)^4&=p(y)^4-q(y)^4\end{aligned}
\]
Die Reihen für $p,q,r$, welche nach ganzen Potenzen von $y^{\frac14}$ fortschreiten und diesen Bedingungen genügen, findet man, so weit man die Entwickelung ausführt von der Form:
\[
\begin{aligned}p(y)&=1+2y+2y^4+2y^9\\&\quad+.\;.+2y^{nn}+.\;.\;.\end{aligned}
\]
\[
\begin{aligned}q(y)&=1-2y+2y^4-2y^9\\&\quad+.\;.\pm2y^{nn}\mp.\;.\;.\end{aligned}
\]
\[
\begin{aligned}r(y)&=2y^{\frac14}+2y^{\frac94}+2y^{\frac{25}4}+2y^{\frac{49}4}\\&\quad+.\;.+2y^{(n+\frac12)^2}+.\;.\;.\end{aligned}
\]

\SourceLeaf{393}%
Dass das hier angedeutete Gesetz für die Bildung der Glieder das allgemein gültige ist, scheint \textsc{Gauss} unter Anderem auch auf folgende Art bewiesen
zu haben. Neben den entsprechenden Gleichungen, welche sich auf Reihen mit
zwei Argumenten beziehen und weiter unten in Art.\,23 und 25 Platz finden werden, hat \textsc{Gauss} sich die Aufzeichnung gemacht:]

\begingroup\itshape
Zur Theorie der Zerlegung der Zahlen in vier Quadrate.

Das Theorem: das Product zweier Summen von vier Quadraten ist selbst eine
Summe von vier Quadraten, wird am einfachsten so dargestellt:

es seien $l,m,\lambda,\mu,\lambda',\mu'$ sechs complexe Zahlen, so dass $\lambda,\lambda'$ und $\mu,\mu'$ sociirt sind.
Durch $\mathrm N$ bezeichne man die Norm. Es ist dann
\[
(\mathrm Nl+\mathrm Nm)(\mathrm N\lambda+\mathrm N\mu)
=\mathrm N(l\lambda+m\mu)+\mathrm N(l\mu'-m\lambda')
\]
\par\endgroup
[und also auch
\[
\begin{aligned}
&\{\mathrm N(n+in_{\prime})+\mathrm N(n_{\prime\prime}+in_{\prime\prime\prime})\}
\\&\quad\{\mathrm N(1-i)+\mathrm N(1+i)\}\\
&\qquad=\mathrm N\{(n+n_{\prime}+n_{\prime\prime}-n_{\prime\prime\prime})
\\&\qquad\quad+i(-n+n_{\prime}+n_{\prime\prime}+n_{\prime\prime\prime})\}\\
&\qquad\quad+\mathrm N\{(n+n_{\prime}-n_{\prime\prime}+n_{\prime\prime\prime})
\\&\qquad\quad-i(+n-n_{\prime}+n_{\prime\prime}+n_{\prime\prime\prime})\}
\end{aligned}
\]
Hieraus lassen sich leicht die beiden folgenden Sätze ableiten, in welchen
verschiedene Darstellungen einer Zahl durch eine Summe von vier Quadratzahlen
sich beziehen auf die verschiedenen Werthensysteme der vier Wurzeln mit Berücksichtigung sowohl der Zeichen als auch der Reihenfolge der Wurzeln, worin
ferner unter den geraden Zahlen auch die Null mit begriffen wird.

Ist das Vierfache einer Zahl von der Form $4k+1$ durch vier ungerade
Quadratzahlen darstellbar, so ist sie selbst halb so oft durch eine ungerade und
drei gerade Quadratzahlen darstellbar, und umgekehrt, ist eine Zahl in der letzten Weise darstellbar, so ist ihr Vierfaches doppelt so oft in der ersten Weise
darstellbar.

Ist das Vierfache einer Zahl von der Form $4k+3$ durch vier ungerade
Quadratzahlen darstellbar, so ist sie selbst halb so oft durch eine gerade und
drei ungerade Quadratzahlen darstellbar und umgekehrt, ist eine Zahl in der letzten Weise darstellbar, so ist ihr Vierfaches doppelt so oft in der ersten Weise
darstellbar.

Mit Hülfe dieser beiden Sätze lässt sich unmittelbar beweisen, dass die
obigen Reihen der Gleichung $r(y)^4=p(y)^4-q(y)^4$ genügen, und ferner durch \SourceLeaf{394}%
die wiederholte Anwendung dieser Relation, dass, wenn man nach dem Schema
\[
p+q=2p'',\quad p-q=2r'',\quad(p'')^4-(r'')^4=(q'')^4
\]
aus den Werthen der Quadrate der beiden ersten Reihen, nemlich $pp$ und $qq$,
als Anfangsglieder die Glieder des Algorithmus eines arithmetisch geometrischen
Mittels für positive gerade Indices bildet, auch
\[
p^{2\mathrm n}=p(y^{2^{2n}}),\quad q^{2\mathrm n}=q(y^{2^{2n}}),\quad r^{2\mathrm n}=r(y^{2^{2n}})
\]
wird. Durch Übergang zu dem Grenzwerthe von $n$ entsteht also nach Art.\,12:
\[
\mathrm M(pp,qq)=1,\qquad
\frac\pi2\frac{\mathrm M(pp,qq)}{\mathrm M(pp,rr)}=-\tfrac12\log y\ ]
\]

\AGMArticle{17}
[Aus den im vorhergehenden Artikel abgeleiteten Eigenschaften der durch
die dort aufgestellten Reihen definirten Functionen $p,q,r$ folgt, dass, wenn
$a,b,c$ drei die Gleichung $aa=bb+cc$ erfüllende Grössen sind, sie in die Form
gesetzt werden können
\[
\begin{gathered}a=\mathrm M(a,b)\cdot(py)^2,\\[5pt] b=\mathrm M(a,b)\cdot(qy)^2,\\[5pt] c=\mathrm M(a,b)\cdot(ry)^2\end{gathered}
\]
und dass dann $\frac{\log y}{\pi}=-\frac{\mathrm M(a,b)}{\mathrm M(a,c)}$ sein muss. Setzt man nun noch
\[
\begin{gathered}a=\mathrm M(a,c)\cdot(pz)^2,\\[5pt] c=\mathrm M(a,c)\cdot(qz)^2,\\[5pt] b=\mathrm M(a,b)\cdot(rz)^2\end{gathered}
\]
so wird
\[
\frac{\log z}{\pi}=-\frac{\mathrm M(a,c)}{\mathrm M(a,b)}=\frac{\pi}{\log y}.
\]
Der durch die Vereinigung dieser beiden Darstellungen sich ergebende Satz
ist von \textsc{Gauss} so ausgesprochen, dass die Functionen
\[
\begin{aligned}\mathfrak P t&=1+2e^{-\pi t}+2e^{-4\pi t}+2e^{-9\pi t}\\&\quad+.\;.+2e^{-nn\pi t}+.\;.\end{aligned}
\]
\[
\begin{aligned}\mathfrak Q t&=1-2e^{-\pi t}+2e^{-4\pi t}-2e^{-9\pi t}\\&\quad+.\;.\pm2e^{-nn\pi t}\mp.\;.\end{aligned}
\]
\[
\begin{aligned}\mathfrak R t&=2e^{-\frac14\pi t}+2e^{-\frac94\pi t}+2e^{-\frac{25}4\pi t}\\&\quad+.\;.+2e^{-(n+\frac12)^2\pi t}+.\;.\end{aligned}
\]
den Gleichungen
